% =========================================================================
% Guía de lectura — presenta las cajas de apoyo y sus patos
% =========================================================================
\chapter*{Guía de lectura}
\addcontentsline{toc}{chapter}{Guía de lectura}

Cada problema parte con su enunciado y su figura, y termina con una caja gris con la
solución desarrollada. A lo largo de las soluciones te vas a encontrar con algunos patos que
nos acompañan con comentarios. Cada pato tiene un trabajo distinto:

\patotip{\textbf{Pato ingeniero.} Da consejos de cómo abordar el problema: qué marco
conviene elegir, qué ecuación usar, en qué fijarse primero.}

\patoojo{\textbf{Pato del letrero.} Aparece cuando hay una trampa: un error típico, una
figura engañosa o un signo que se pierde fácil.}

\patotruco{Atajos que ahorran álgebra. No son obligatorios, pero una vez que los
entiendes no los vas a querer soltar.}

\begin{resumen}[Pato profe: resúmenes]
Al final de un tema o de un método, el pato profe deja lo importante en una caja azul:
las fórmulas que hay que tener a mano y los pasos de la receta.
\end{resumen}

\begin{solucion}{Nombre del problema}
Las soluciones van en una sola caja como esta. Si estás usando el libro para practicar,
intenta el problema antes de mirarla.
\respuesta{Los resultados finales quedan destacados así.}
\end{solucion}
\clearpage
